\documentclass[12pt]{article}

\usepackage[utf8]{inputenc}
\usepackage{latexsym,amsfonts,amssymb,amsthm,amsmath}
\usepackage{tikz}
\usepackage{stanli}
\usetikzlibrary{shapes,calc,through,3d,patterns}
\usetikzlibrary{arrows.meta,decorations.pathreplacing,angles,quotes}
\usepackage{multicol}
\usepackage{geometry}
\usepackage{mathtools}
\usepackage{siunitx}
\usepackage{booktabs}

\setlength{\parindent}{0in}
\setlength{\oddsidemargin}{-0.25in}
\setlength{\textwidth}{7in}
\setlength{\textheight}{9.3in}
\setlength{\topmargin}{-1in}
\setlength{\headheight}{18pt}

\newcommand\blfootnote[1]{%
    \begingroup
    \renewcommand\thefootnote{}\footnote{#1}%
    \addtocounter{footnote}{-1}%
    \endgroup
}

\definecolor{darkred}{rgb}{0.6,0,0}

\newcommand{\bolddelta}{\boldsymbol\delta}
\newcommand{\boldxi}{\boldsymbol\xi}

\newcommand{\T}{^{T}}
\newcommand{\bK}{\overline{\mathbf K}}
\newcommand{\bM}{\overline{\mathbf M}}
\newcommand{\vct}[1]{\mathbf{#1}}
\newcommand{\colv}[1]{\begin{bmatrix}#1\end{bmatrix}}

\begin{document}

\noindent ME473 - Dynamic finite element analysis of structures\hfill EPFL\\
Stefano Burzio \hfill 2025\\
\vspace{0.3cm}

\hspace*{\fill}\textbf{\Large Problem set 10}\hspace*{\fill}

\hrulefill

\subsection*{Problem 1}

Let $\vct K$ and $\vct M$ be a symmetric positive definite matrices of order $q$. Pick $i<j$ and use the $2\times 2$ blocks
\[
  \vct K_{(i,j)}=\colv{k_{ii}&k_{ij}\\k_{ij}&k_{jj}},\qquad
  \vct M_{(i,j)}=\colv{m_{ii}&m_{ij}\\m_{ij}&m_{jj}} .
\]
as each Jacobi transformation acts only on coordinates $i$ and $j$.
\begin{enumerate}
  \item[a)] One pivot in the standard Jacobi problem. Take $\vct M_{(i,j)}=\vct I$ and the Jacobi transformation matrix to be the rotation
        \[
          \vct J(\theta)=
          \colv{\cos\theta&-\sin\theta\\ \sin\theta&\cos\theta}.
        \]
        Compute the components of the rotated matrices
        $$\vct K_{(i,j)}'=\vct J\T\vct K_{(i,j)}\vct J,\quad\mbox{and}\quad  \vct M_{(i,j)}'=\vct J\T\vct M_{(i,j)}\vct J$$
        and verify that, if the angle is chosen as
        \[
          \tan(2\theta)=\frac{2k_{ij}}{k_{ii}-k_{jj}},\qquad |\theta|\leq\pi/4 .
        \]
        then the matrix $\vct K_p'$ is diagonal and $\vct M_p'$ is the identity matrix. Explain the choice of angle when $k_{ij}=0$ or $k_{ii}=k_{jj}$.
  \item[b)] One pivot pair in the generalized problem. Remove the assumption that $\vct M_{(i,j)}=\vct I$, and modify the definition of Jacobi transformation matrix to
        $$\vct J(\alpha,\gamma)=\colv{1&\alpha\\\gamma&1}.$$
        Show that simultaneous elimination of the stiffness and mass off-diagonal entries in $\vct K_{(i,j)}'=\vct J\T\vct K_{(i,j)}\vct J$ and $\vct M_{(i,j)}'=\vct J\T\vct M_{(i,j)}\vct J$ requires
        \begin{align*}
          k_{ii}\alpha+k_{jj}\gamma+k_{ij}(1+\alpha\gamma) & =0, \\
          m_{ii}\alpha+m_{jj}\gamma+m_{ij}(1+\alpha\gamma) & =0.
        \end{align*}
  \item[c)] Find $\alpha$ and $\gamma$ by introduce the quantities
        \[
          \bar k_{ii}=k_{ii}m_{ij}-m_{ii}k_{ij},\quad
          \bar k_{jj}=k_{jj}m_{ij}-m_{jj}k_{ij},\quad
          \bar k=\tfrac12(k_{ii}m_{jj}-k_{jj}m_{ii}).
        \]
        Eliminate $1+\alpha\gamma$ to find a relation between the two unknowns.
        Write $\alpha=\bar k_{jj}/x$, $\gamma=-\bar k_{ii}/x$ and derive the
        expression of $x$. Which root has the larger magnitude, and why
        is it preferable when evaluating the formula?
  \item[d)] For $\vct M_p=\vct I$, show that the argument developed in points b) and c) allows to recovers $\vct J(\theta)$ as in part a).
\end{enumerate}


\subsection*{Problem 2}
Develop a \texttt{MATLAB} script to compute the first two natural frequencies and the corresponding mode shapes associated with the longitudinal vibrations of a uniform free-free bar, using the subspace iteration method. The structure is characterized by the following physical parameters:
\begin{multicols}{2}
  \begin{itemize}
    \item length: $\ell=0.5$ m,
    \item cross-sectional area: $A=4\,10^{-4}$ m$^2$,
    \item Young's modulus: $E=210$ GPa,
    \item mass density: $\rho=7850$ kg/m$^3$.
  \end{itemize}
\end{multicols}
The bar is modeled using a single quadratic finite element with three equally spaced nodes. The structure stiffness and lumped mass\footnote{The use of a lumped mass matrix is often justified by its ability to simplify computations, improve numerical stability in explicit time integration schemes, and avoid non-physical oscillations associated with consistent mass formulations in low-order elements.} matrices are given by:
\[
  \mathbf{K} = \frac{EA}{3\ell}
  \begin{bmatrix}
    7  & -8 & 1  \\
    -8 & 16 & -8 \\
    1  & -8 & 7
  \end{bmatrix}\qquad\text{and}\qquad
  \mathbf{M} = \frac{\rho A \ell}{6} \text{diag}(1, 4, 1).
\]

\begin{enumerate}
  \item[1)] Formulate the generalized eigenvalue problem for the system, and define the stiffness and mass matrices explicitly.

  \item[2)] Since the system exhibits a rigid-body mode with zero frequency, a spectral shift is applied to ensure convergence. The modified eigenvalue problem is:
        \[
          (\mathbf{K} + \sigma \mathbf{M}) \mathbf{p} = (\lambda + \sigma) \mathbf{M} \mathbf{p}
        \]
        where the shift parameter is chosen (to simplify the computation) as:
        \[
          \sigma = 2\frac{E}{\rho \ell^2}.
        \]
        Define the shifted stiffness matrix as $\mathbf{K}_\sigma = \mathbf{K} + \sigma \mathbf{M}$, yielding the equivalent eigenvalue problem:
        \[
          \mathbf{K}_\sigma \mathbf{p} = \lambda_\sigma \mathbf{M} \mathbf{p}, \quad \text{with} \quad \lambda_\sigma = \lambda + \sigma.
        \]

  \item[3)] Implement three iterations of the subspace iteration algorithm to approximate the two lowest eigenpairs $(\lambda_i, \mathbf{p}_i)$. For simplicity use the function \texttt{eig} to solve the reduced Rayleigh's minimization problem. Use the following initial guess:
        \[
          \mathbf{P}_{(0)} = \frac{1}{\sqrt{\rho A \ell}} \begin{bmatrix}
            1 & 0 \\
            0 & 1 \\
            0 & 0
          \end{bmatrix}.
        \]

  \item[4)] Compare the numerical approximation with the exact (shifted) eigenvalues and corresponding eigenvectors (defined up to a multiplicative constant):
        \[
          \lambda_{\sigma,1} = 2\frac{E}{\rho \ell^2}\qquad\text{and}\qquad \lambda_{\sigma,2} = 14\frac{E}{\rho \ell^2},
        \]
        \[
          \mathbf{p}_1 = \begin{bmatrix} -1 \\ -1 \\ -1 \end{bmatrix}\qquad\text{and}\qquad
          \mathbf{p}_2 = \begin{bmatrix} -1 \\ 0 \\ 1 \end{bmatrix}
        \]
        and evaluate the accuracy by computing:
        \begin{itemize}
          \item The relative error between approximated and exact eigenvalues.
          \item The error in the $\mathbf{M}$-norm between approximated and exact eigenvectors.
        \end{itemize}

        %\item Investigate the sensitivity of convergence to the choice of initial vectors. Repeat the computation using the alternative initial matrix:
        %   \[
        %    \mathbf{P}_{(0)} = \frac{1}{\sqrt{\rho A \ell}} \begin{bmatrix}
        %      1 & 0 \\
        %      4 & 1 \\
        %      1 & 0
        %    \end{bmatrix}
        %  \]
        %  and discuss the observed changes in convergence behavior.
\end{enumerate}


\end{document}
